\section{Concluding remarks}\label{Section5}
In this study, a phase-field fracture model for ferroelectric materials that incorporates the flexoelectric effect is established. The model uses isogeometric analysis method with PHT-splines to realize the local refinement near the crack. The crack propagation process in ferroelectric materials with electrically impermeable crack conditions is simulated. The results indicate that the crack propagation in ferroelectric materials depends on the polarization direction with the consideration of flexoelectric effect. When the crack propagation direction aligns with the polarization direction, the crack extends the fastest. In contrast, the crack extension is decelerated when the propagation direction opposes the polarization direction. When the initial polarization is perpendicular to the crack, the flexoelectric effect can induce a deflection in the crack path, which is opposite to the initial polarization direction. The underlying mechanism can be attributed to the eigenstrain induced by the flexoelectric effect.

The numerical results are consistent with previous experimental observations  and  provide a phenomenological explanation for the anisotropic effects during fracture in ferroelectric materials. From this work, some further researches can also be studied. For example, this study used an electrically impermeable crack. The impact of different electrical boundary conditions on crack propagation still needs investigation. Additionally, a two-dimensional model was used to simplify calculations. However, for materials like PMN-PT, a two-dimensional model cannot accurately describe their domain structures. Therefore, a three-dimensional model is still needed to study the fracture process.